\section{Introduction}
\label{sec:intro}

Video anomaly detection (VAD) systems are increasingly deployed at the edge, on cameras, robots, and mobile devices, where compute, memory, and power are all constrained and predictions must be produced online, as frames arrive, rather than after the fact. IEEE Transactions on Pattern Analysis and Machine Intelligence has an active recent literature on this problem, from a taxonomy and evaluation survey of single-scene methods~\cite{jones2022survey} to the journal extension of the future-frame-prediction baseline~\cite{luo2022futureframe}, a scene-aware latent-space prediction and anticipation model~\cite{cao2025scenedependent}, and a self-supervised detector based on learned neural transformations~\cite{qiu2025neutral}. None of these, however, is designed around the specific constraint that motivates this paper: strict causality, meaning a prediction at time $t$ must depend only on frames $1,\dots,t$ with bounded per-step compute and memory, combined with real hardware feasibility, meaning latency and throughput are measured rather than asserted on the class of device the method targets.

State-space models (SSMs), and Mamba~\cite{gu2023mamba} in particular, have recently been proposed for VAD on efficiency grounds. VADMamba~\cite{vadmamba2025} and its follow-up VADMamba++~\cite{vadmambaplusplus2026} combine a Mamba-based temporal block with vector-quantized frame prediction and optical-flow reconstruction, and Wave-MambaAD~\cite{wavemambaad2025} adds a wavelet-driven SSM for unified subtle and large-scale anomaly detection. These works establish that SSMs are a good fit for VAD's efficiency requirements, but they still process the input in clips or windows rather than truly online per frame, report no theoretical relationship between the model's temporal memory and how quickly it can detect an anomaly, and benchmark efficiency in terms of GPU throughput rather than measuring the target-deployment hardware directly. This paper closes those three gaps together, rather than any one in isolation, because they compound: a model that is only softly causal cannot make a hard latency guarantee, and a latency number without a theory of why the model reacts when it does is an empirical curiosity rather than a design tool.

This paper makes three contributions. First, a strictly causal streaming formulation of SSM-based VAD, with a concrete architecture (Section~\ref{sec:method}) in which every operation is expressible as a per-frame $O(1)$ update. Second, a theoretical analysis (Section~\ref{sec:theory}) relating the recurrence's decay spectrum to detection delay and minimum-detectable anomaly duration, validated empirically. Third, an end-to-end evaluation including real on-device latency, memory, and throughput measurements on consumer edge silicon (Section~\ref{sec:experiments}), alongside standard accuracy metrics and an ablation study that we report in full, including a result that complicates rather than confirms our own design hypothesis.
